\chapter{C++的多线程编程}

我们刚学编程的时候，代码大多数时候是一个线程跑到尾。那对于一些大型计算任务，比如大型游戏，一个线程从头跑到尾，效率太低。现代CPU动辄几十个核，难道一核有难、八核围观？这显然是不合适的。还有另一些情形：假如现在有两个任务，一个是要从网上下载某个超大型资源，另一个是要处理这些资源。朴素的想法是先把所有东西全都拉下来然后再一起处理，但这样就会造成资源的浪费：下载的时候CPU在睡觉，不能把这段时间给利用起来。

\section{前置知识：程序执行原理}

在深入学习多线程编程之前，我们需要了解一些基本的程序执行原理：操作系统是如何调度程序的，CPU是如何执行程序的。

\subsection{进程}

操作系统是一个资源管理器，它管理着计算机的各种资源，包括CPU、内存、磁盘、网络等。

当一个程序执行起来的时候，操作系统会为它分配一个独立的内存空间，并为它创建一个\textbf{进程}，即“进行中的程序”。多数程序仅由一个进程构成，但也有一些程序会创建多个进程来完成不同的任务。每个进程都有自己的内存空间和资源，它们之间是相互隔离的。

进程是操作系统调度资源的基本单位。当一个进程被创建时，操作系统会为该进程分配一个唯一的“进程标识符”（Process ID, PID），并为其分配内存空间和其他资源。进程中的代码会被加载到内存中，并开始执行。这个过程的资源开销是比较大的。由于不同进程的内存空间相互隔离，因此要让进程之间通信，必须经过系统调用，这个开销也是较大的。

为了允许计算机执行多个线程，操作系统会进行时间片轮转调度。时间片是CPU分配给每个进程的时间段，操作系统会在不同进程之间切换，让它们轮流使用CPU。这样看起来就像是多个进程同时在运行，但实际上在任意时刻，CPU只会执行一个进程的指令。

\subsection{线程}

线程是CPU调度的基本单位。可以理解为一个单一的计算任务或执行路径。一个进程可以包含多个线程，这些线程共享进程的内存空间和资源。线程之间的切换开销比进程之间的切换开销小得多。这样，一个复杂的计算任务（进程）就可以被拆分成许多小的线程，并让它们同时运行在不同的CPU核心上，从而提升整体的计算效率。打个比方，进程是“项目组”，而线程则是“员工”；领取资源是以“项目组”的名义领取的，而具体的工作则是由“员工”来完成的。通过合理地分配任务，不同的员工同时进行不同的工作，就能提高整个项目组的工作效率。

\section{基本多线程编程}

\subsection{从异步说开去}

C++中的异步本质上也是多线程编程的一种形式，即 \verb|std::async|。异步是C++中最方便的启动多线程任务的一种方式。它可以让你在后台执行一个函数，并在需要的时候获取其结果，而不需要显式管理线程的创建和销毁。

\begin{minted}{cpp}
#include <future>
#include <thread>
#include <chrono>
#include <print>

int compute(){
    // 假装这是一个耗时的计算任务，比如两个大有序数组的归并
    // 这里只是让这个线程睡 10 秒，模拟计算的耗时
    std::println("Computing...");
    std::this_thread::sleep_for(std::chrono::seconds(10));
    return 42;
}

int main() {
    // 异步的启动 compute() 函数，然后啥都不用管，结果在需要的时候再去获取
    std::future<int> result = std::async(std::launch::async, compute);

    // 该干嘛干嘛
    std::println("Doing other things...");

    // 别的干完了后回来获取结果
    // 如果结果没计算完，这里会在这里等着
    int value = result.get();
    std::println("Result: {}", value);

    return 0;
}
\end{minted}

这就是最基本的多线程编程。在这个例子中，\verb|std::future| 指的是“未来的值”，它是一个占位符，表示某个值将在未来的某个时间点被计算出来。通过调用 \verb|get()| 方法，我们可以获取这个值。如果计算还没有完成，\verb|get()| 会阻塞当前线程，直到计算完成并返回结果。

\mintinline{cpp}{std::async} 是一个函数模板，其签名类似于：
\begin{minted}{cpp}
template<class F, class... Args>
std::future<std::invoke_result_t<F, Args...>> async(F&& f, Args&&... args);
\end{minted}
它接受一个可调用对象（函数、函数指针、lambda表达式等）和一组参数，并返回一个 \verb|std::future| 对象。这个对象可以用来获取异步计算的结果。而在这里的\verb|std::launch::async| 是一个枚举值，表示异步任务应该在新的线程中执行。相应的，\verb|std::launch::deferred| 表示异步任务应该在调用 \verb|get()| 时才执行，即延迟执行。

\subsection{发生了什么？}

让我们看看刚刚这个程序在做什么。

用GDB来调试这个程序，得到的结果如下：
\begin{minted}{text}
Breakpoint 1, main () at test.cc:14
14      int main()
(gdb) n
16              std::future<int> result = std::async(std::launch::async, compute);
(gdb) n
[New Thread 0x7ffff77ff6c0 (LWP 228816)]

Thread 1 "test" hit Breakpoint 3, main () at test.cc:19
19              std::println("Doing other things...");
(gdb) n
[Switching to Thread 0x7ffff77ff6c0 (LWP 228816)]

Thread 2 "test" hit Breakpoint 2, compute () at test.cc:6
6       int compute()
(gdb) n
Doing other things...
[Switching to Thread 0x7ffff7f4b780 (LWP 228789)]

Thread 1 "test" hit Breakpoint 4, main () at test.cc:23
23              int value = result.get();
(gdb) n
Computing...
[Thread 0x7ffff77ff6c0 (LWP 228816) exited]
24              std::println("Result: {}", value);
(gdb) n
Result: 42
26              return 0;
\end{minted}

根据GDB的调试信息，我们可以看到：当 \mintinline{cpp}{std::async(std::launch::async, compute)} 被调用时，C++会创建一个新的线程来执行 \verb|compute()| 函数。这个新线程会在后台运行，而主线程继续执行后续的代码。

所以说，如果显式的声明一个线程，代码就会这样：
\begin{minted}{cpp}
#include <future>
#include <thread>
#include <chrono>
#include <print>
#include <promise>

int compute(){
    std::println("Computing...");
    std::this_thread::sleep_for(std::chrono::seconds(10));
    return 42;
}

int main(){
    std::promise<int> promise;
    std::future<int> result = promise.get_future();

    // 显式创建一个线程来执行 compute() 函数
    std::thread t([&promise](){
        int value = compute();
        promise.set_value(value);
    });

    std::println("Doing other things...");

    int value = result.get();
    std::println("Result: {}", value);

    t.join(); // 等待线程结束

    return 0;
}
\end{minted}
可以看到，新的线程发起了一个“承诺”，承诺它将在未来的时间点提供一个数值。然后，在这个承诺下，线程自己去跑 \verb|compute()| 函数，计算完毕后将结果放入承诺中。主线程则继续执行其他任务，直到需要结果时才去获取它。

\subsection{线程生命周期管理}

在多线程编程中，线程的生命周期管理是一个重要的问题。整个进程结束时，所有子线程必须结束，随后主要线程才能结束，否则会引发未定义行为。

通常，我们会使用 \verb|std::thread| 的 \verb|join()| 方法来等待线程完成，或者使用 \verb|detach()| 方法让线程独立运行。

\verb|.join()| 方法的意思是：阻塞当前的线程，直到本线程完成执行。完成执行后，线程会被销毁，资源会被释放。使用 \verb|.join()| 可以确保线程在程序结束前完成任务，避免出现未定义行为。C++的RAII在这里也发挥了空间，C++20的 \verb|std::jthread| 会在该线程析构时自动调用 \verb|.join()| 方法，确保线程在程序结束前完成任务。此外，该类还支持 \verb|stop_token|，可以在需要时优雅地停止线程。用 \verb|std::jthread| 重写上述程序如下：
\begin{minted}{cpp}
#include <jthread>
#include <future>
#include <print>

int compute(){
    std::println("Computing...");
    std::this_thread::sleep_for(std::chrono::seconds(10));
    return 42;
}

int main(){
    std::promise<int> promise;
    std::future<int> result = promise.get_future();

    // 使用 jthread 来管理线程生命周期
    std::jthread t([&promise](){
        int value = compute();
        promise.set_value(value);
    });

    std::println("Doing other things...");

    int value = result.get();
    std::println("Result: {}", value);

    // jthread 会在析构时自动 join
    return 0;
}
\end{minted}

\verb|.detach()| 方法的意思是：让线程脱离主线程的控制，独立运行。这样，主线程不会等待该线程完成，而是继续执行后续代码。使用 \verb|.detach()| 需要特别小心，因为如果主线程在子线程完成之前就结束了，那么子线程可能会在后台继续运行，导致资源泄露或程序崩溃。

有了这个基础，我们就可以在C++中进行多线程编程了。

\begin{example}{}{}

请对一个一千万个整数的数组进行排序。

\tcblower

字越少，事越大。一千万个整数就已经不是一般单线程程序能应付得了的了。

怎样去写这个多线程编程呢？我们想到了归并排序，它天生适合多线程编程。我们可以将数组分成两半，分别在两个线程中进行排序，然后再将两个有序数组合并成一个有序数组。这个过程可以递归地进行下去，直到每个线程只需要排序少数元素为止（1000个元素以内的排序可以直接使用 \verb|std::sort()| 来完成）。这样，我们就可以充分利用多核CPU的计算能力，大大提高排序的效率。代码实现大致如下：

\begin{minted}{cpp}
#include <algorithm>
#include <future>
#include <print>
#include <random>
#include <vector>

std::vector<int> parallel_merge_sort(std::vector<int>& arr) {
	if (arr.size() <= 1000) {
		std::sort(arr.begin(), arr.end());
		return arr;  // 返回原对象（RVO/移动）
	}

	size_t mid = arr.size() / 2;
	std::vector<int> left(arr.begin(), arr.begin() + mid);
	std::vector<int> right(arr.begin() + mid, arr.end());

	// 左半异步，右半同步，防止线程过多
	auto left_future = std::async(std::launch::async, parallel_merge_sort, std::ref(left));
	auto right_sorted = parallel_merge_sort(right);   // 同步排序

	auto left_sorted = left_future.get();  // 等待左半完成

	std::vector<int> sorted(left_sorted.size() + right_sorted.size());
	std::merge(left_sorted.begin(), left_sorted.end(),
			   right_sorted.begin(), right_sorted.end(),
			   sorted.begin());
	return sorted;
}

int main() {
	std::vector<int> arr(10000000); // 一千万个元素
	std::ranges::iota(arr, 0);       // 填充数组
	std::ranges::shuffle(arr, std::mt19937{std::random_device{}()}); // 打乱数组
	std::println("Array is shuffled!");

	std::vector<int> sorted_arr = parallel_merge_sort(arr);

	// 检查一下
	if (std::ranges::is_sorted(sorted_arr)) {
		std::println("Array is sorted!");
	} else {
		std::println("Array is NOT sorted!");
	}
	return 0;
}
\end{minted}

当然也可以投机取巧：
\begin{minted}{cpp}
#include <algorithm>
#include <print>
#include <vector>
#include <execution>
#include <random>

int main() {
    std::vector<int> arr(10000000); // 一千万个元素
    std::ranges::iota(arr, 0);       // 填充数组
    std::ranges::shuffle(arr, std::mt19937{std::random_device{}()}); // 打乱数组
    std::println("Array is shuffled!");

    // 使用并行执行策略进行排序，见上一章
    std::sort(std::execution::par, arr.begin(), arr.end());

    // 检查一下
    if (std::ranges::is_sorted(arr)) {
        std::println("Array is sorted!");
    } else {
        std::println("Array is NOT sorted!");
    }
    return 0;
}
\end{minted}
这种写法确实是多线程。实际情形中这么写也不是不行。

\end{example}

上述写法（指手搓的多线程实现）在数组长度为一千万左右的时候表现尚且良好。但当数组长度到达一亿甚至更高的时候这个方法失效。理由是：线程太多，操作系统并不能很好的调度如此海量的线程，导致线程之间的切换开销过大；其次，虽然创建和销毁线程的开销远小于进程，但也不是没有开销，创建和销毁太多线程也会导致性能下降。最后，线程的数量过多会导致内存消耗过大，甚至可能导致系统崩溃。

那么，怎么去解决这些问题呢？

\section{线程池}

既然创建线程的开销不小，那么我们为什么还要频繁创建、销毁线程呢？如果预先创造好一组线程，并让它们睡觉、等待任务的到来；当有任务时，便唤醒一些线程，让他们去干活，干完活继续睡觉，而不是被销毁；如果任务到来时，线程池中的线程都在干活，那么就让任务排队等待，等有线程空闲了再去干活。这样就可以大大减少线程的创建和销毁开销，还能控制线程的总数。而按照这个思路写出来的东西，就是\textbf{线程池}。

\subsection{线程池的基本原理和实现}

那么问题是：怎么给这群线程派任务？任务来了，谁去干活？怎么保证这个任务仅唤醒一个线程去干活？怎么保证唤醒多个线程时，这些线程不会抢着做同一件事？这些问题都需要我们去解决。

在线程池的实现中，我们通常会让所有的线程都在 \verb|while(true)| 这个死循环中消极等待。任务被包装成一个个 \verb|std::function<void()>| 对象，放入一个任务队列中，并由主要线程通知一个线程起来干活。为了进一步保证仅有一个线程在干这个活（因为有时候线程确实可能在没有收到通知的时候被意外唤醒），我们要给这个任务队列加上锁，保证一次只有一个线程能访问之（即“获得锁”）；当其他进程获得锁时，队列中已经没有那个任务了，此时不会重复执行，而是继续睡觉。


因此，一个简单而功能完整的固定大小线程池的实现大致如下：
\begin{minted}{cpp}
#include <vector>
#include <thread>
#include <queue>
#include <mutex>
#include <condition_variable>
#include <functional>
#include <atomic>
#include <future>  
#include <iostream>
#include <stdexcept>

class ThreadPool {
    std::vector<std::thread> workers; // 工作线程们
    std::queue<std::function<void()>> tasks; // 任务队列
    std::mutex queue_mutex; // 任务队列的锁
    std::condition_variable condition; // 条件变量，用于通知线程有新任务
    std::atomic<bool> stop; // 停止标志，用于关闭线程池

    void workerLoop() { // 睡觉循环
        while (true) {
            std::function<void()> task;
            {
                // 进块，获取“RAII”的锁
                std::unique_lock<std::mutex> lock(queue_mutex);

                // 等待条件变量，直到有任务或者线程池被停止
                condition.wait(lock, [this] { return stop || !tasks.empty(); });
                if (stop && tasks.empty())
                    return; // 如果停止并且任务队列为空，退出线程
                
                // 获取任务
                task = std::move(tasks.front());
                tasks.pop();
            } // 出块，自动释放锁
            if (task) { task(); }
        }
    }

public:
    // 构造，启动这些线程，默认为 CPU 核心数
    ThreadPool(size_t threads = std::thread::hardware_concurrency()) : stop(false) {
        for (std::size_t i = 0; i < threads; ++i)
            workers.emplace_back(&ThreadPool::workerLoop, this);
    }

    // 析构，停止线程池，此时要保证所有线程都退出
    ~ThreadPool() {
        {
            std::unique_lock<std::mutex> lock(queue_mutex); // 给任务队列上锁，保证线程安全
            stop = true; // 设置停止标志
        }
        // 通知所有线程，唤醒它们
        condition.notify_all();
        // 等待所有线程退出
        for (std::thread &worker : workers)
            if (worker.joinable())
                worker.join();
    }

    // 禁止拷贝和赋值线程池
    ThreadPool(const ThreadPool&) = delete;
    ThreadPool& operator=(const ThreadPool&) = delete;

    // 干活（void）
    void enqueue(std::function<void()> task) {
        {
            std::unique_lock<std::mutex> lock(queue_mutex); // 给任务队列上锁，保证线程安全
            if (stop) throw std::runtime_error("enqueue on stopped ThreadPool");
            tasks.emplace(std::move(task)); // 将任务放入队列
        }
        condition.notify_one(); // 通知一个线程去干活
    }

    // 干活（有返回值），需要用完美转发和 std::packaged_task 来实现
    template<class F, class... Args>
    auto enqueue(F&& f, Args&&... args) -> std::future<std::invoke_result_t<F, Args...>> {
        using return_type = std::invoke_result_t<F, Args...>; // 把返回值类型写的简单些

        // 包装任务。这里必须用共享所有权的智能指针而不能用独占所有权的智能指针，因为任务队列中可能有多个线程在等待这个任务的完成。
        auto task = std::make_shared<std::packaged_task<return_type()>>(
            [f = std::forward<F>(f), ...args = std::forward<Args>(args)]() mutable { return f(args...); }
        );
        // 新写法是用匿名函数，基本已经取代了 std::bind。
        
        std::future<return_type> res = task->get_future(); // 获取任务的 future
        {
            std::unique_lock<std::mutex> lock(queue_mutex); // 给任务队列上锁，保证线程安全
            if (stop) throw std::runtime_error("enqueue on stopped ThreadPool");
            tasks.emplace([task]() { (*task)(); }); // 将任务放入队列并包装成 void() 类型
        }
        condition.notify_one(); // 通知一个线程去干活
        return res; // 返回 future
    }
};

// 以下是测试 
int addsleep(int a, int b) { // 一个睡1秒的加法，用于模拟耗时计算
    std::this_thread::sleep_for(std::chrono::seconds(1));
    return a + b;
}

int main(){
    ThreadPool pool(4); // 创建一个线程池，大小为4
    std::vector<std::future<int>> results; // 用于存储结果的 future

    for (int i = 0; i < 8; ++i) {
        results.emplace_back(pool.enqueue(addsleep, i, i)); // 提交任务到线程池
    }

    for (auto && result : results) {
        std::cout << result.get() << ' '; // 获取结果并打印
    }
    std::cout << std::endl;

    return 0;
}
\end{minted}
如果正常，上述线程池的测试结果是：停1秒后，输出了0到14之间的所有偶数。

\subsection{C++线程库的同步原语}

现在我们将逐个解释新出现的这些东西：\verb|std::mutex|、\verb|std::condition_variable|、\verb|std::unique_lock|、\verb|std::atomic|、\verb|std::packaged_task|。

以及上文没出现过的四个重要同步原语：\verb|std::shared_mutex|、\verb|std::barrier|、\verb|std::latch|、\verb|semaphore|。

\paragraph{互斥量和锁} 互斥量（\verb|std::mutex|）是一个同步原语，顾名思义就是“互相排斥”的量，主要作用就是保证一个共享的数据（在上述线程池中是任务队列）在同一时刻只能被一个线程访问。

可以这样理解：一个房间就是一个共享数据，一般不会把整个房间（整个共享数据）都设成锁，而是用一个门钥匙来控制谁能进去。一个线程要访问共享数据，就必须先获取这个钥匙（即获取互斥量的所有权），访问完毕后再把钥匙还回去（即释放互斥量的所有权）。因为钥匙只有一个，所以其他线程也想访问这个共享数据时就必须等待，直到钥匙被归还。

普通的锁，类似 \verb|new-delete|，需要手动去获取和释放：
\begin{minted}{cpp}
std::mutex mtx; // 互斥量
mtx.lock(); // 获取互斥量的所有权
// 访问共享数据
mtx.unlock(); // 释放互斥量的所有权
\end{minted}
但是如果在访问共享数据的过程中抛出了异常，那么 \verb|unlock()| 就不会被调用，从而导致互斥量的所有权无法释放，其他线程就无法获取互斥量的所有权，从而导致死锁。为了解决这个问题，C++提供了许多RAII的锁。

最轻的RAII锁是 \verb|std::lock_guard|，它在构造时获取互斥量的所有权，在析构时释放互斥量的所有权。它的使用方式如下：
\begin{minted}{cpp}
std::mutex mtx; // 互斥量
{
    std::lock_guard<std::mutex> lock(mtx); // 获取互斥量
    // 访问共享数据，此时被一直锁住，其他线程无法访问共享数据
} // 离开作用域，lock被析构，释放互斥量
\end{minted}
\verb|std::lock_guard| 不能复制也不能移动，而且其生命周期和作用域高度绑定，不能手动解锁加锁，而且被设计为仅能用于互斥量。在语义明确且操作简单

而更灵活的RAII锁是 \verb|std::unique_lock|，它可以手动加锁解锁，而且可以移动，但不能复制。它的使用方式如下：
\begin{minted}{cpp}
std::mutex mtx; // 互斥量
{
    std::unique_lock<std::mutex> lock(mtx); // 获取互斥量
    // 访问共享数据，此时其他线程无法访问共享数据
    lock.unlock(); // 手动解锁
    // 纯计算，不访问共享数据，此时其他线程可以访问共享数据
    lock.lock(); // 手动上锁
    // 访问共享数据，此时其他线程无法访问共享数据
} // 离开作用域，lock被析构，释放互斥量
\end{minted}
它允许手动解锁加锁、延迟加锁 \verb|std::defer_lock|、尝试加锁 \verb|.try_lock| 等操作。另外，要锁条件变量的话，必须用它。

\paragraph{共享互斥量和读写锁} 共享互斥量（\verb|std::shared_mutex|）是一个同步原语，它允许多个线程同时读取共享数据，但在写入共享数据时，必须保证没有其他线程在读取或写入。它的使用方式如下：
\begin{minted}{cpp}
std::shared_mutex smtx; // 共享互斥量
// 线程A（读线程）
{
    std::shared_lock<std::shared_mutex> lock(smtx); // 获取共享互斥量的共享所有权
    // 读取共享数据，此时其他线程也可以读取共享数据，但不能写入共享数据
} // 离开作用域，lock被析构，释放共享互斥量的共享所有权
// 线程B（写线程）
{
    std::unique_lock<std::shared_mutex> lock(smtx); // 获取共享互斥量的独占所有权
    // 写入共享数据，此时其他线程不能读取或写入共享数据
} // 离开作用域，lock被析构，释放共享互斥量的独占所有权
\end{minted}
共享互斥量的使用场景是读多写少的情况，比如缓存、配置文件等。它可以提高并发性能，但也可能导致写线程饥饿（即写线程一直无法获取独占所有权）。

一般情况下，还是建议使用普通的互斥量和锁，因为它们的语义更简单，容易理解和使用。

\paragraph{条件变量} 条件变量（\verb|std::condition_variable|）是一个同步原语，它允许线程在某个条件不满足时进入睡眠状态，避免空转CPU；当条件满足时，其他线程可以通知它们醒来继续执行。

条件变量有三个主要的用法：睡觉、唤醒一个线程、唤醒所有线程。条件变量的使用方式如下：
\begin{minted}{cpp}
std::condition_variable cv; // 条件变量
std::mutex mtx; // 互斥量
// 线程A
{
    std::unique_lock<std::mutex> lock(mtx); // 获取互斥量
    cv.wait(lock, []{ return condition; }); // 如果条件不满足，进入睡眠状态，直到被唤醒
    // 条件满足，继续执行
} // 离开作用域，lock被析构，释放互斥量
// 主要线程
{
    std::unique_lock<std::mutex> lock(mtx); // 获取互斥量
    // 修改条件
} // 离开作用域，lock被析构，释放互斥量
cv.notify_one(); // 唤醒一个线程
// 或者 cv.notify_all(); // 唤醒所有线程
\end{minted}
\textbf{重要}：用条件变量让线程睡觉时，必须用带谓词的 \verb|.wait|，或把 \verb|.wait| 放在 \verb|while(!condition)| 循环中。原因有二：
\begin{itemize}
    \item 虚假唤醒：操作系统可能无故唤醒线程，这不是bug，但可能导致错误；
    \item 多唤醒竞争：可能唤醒多个线程或错误的线程，也有可能在停止场景下被唤醒但没有任务可做。
\end{itemize}

另外，能唤醒一个线程就尽量不要唤醒所有线程，因为唤醒所有线程会引发“惊群效应”，虽然通过循环检查和队列 \verb|pop| 能保证线程安全（即只有一个线程能拿到任务），但其他线程被唤醒后又要进入睡眠状态，这样会造成CPU资源的浪费。

最后，最佳的实践是先解锁再通知，因为如果先通知再解锁，可能会导致被唤醒的线程立即争抢锁（但此时锁在主线程手里），从而导致被唤醒的线程被阻塞，影响程序性能。

\paragraph{原子变量} 原子变量（\verb|std::atomic|）是一个同步原语，它允许多个线程同时访问同一个变量而不会发生数据竞争。这是一个不需要加锁也能保证线程安全的变量，但只保证单一的读写操作原子性，不能保证复合操作的原子性。原子变量的使用方式如下：
\begin{minted}{cpp}
std::atomic<int> counter(0); // 原子变量
// 线程A
counter++; // 原子操作，线程安全
if (counter.load() > 10) { // 原子操作，线程安全
    // do something
}
// 线程B
counter--; // 原子操作，线程安全
counter = -counter; // 事务性不安全，数据可能被覆盖
\end{minted}
为了保证最后这个事务性操作的逻辑安全（编译上肯定没问题，因为所有操作都是原子的，只是操作复合了，可能在中途 \verb|counter| 就被其他线程修改了），我们必须使用CAS循环（Compare-And-Swap，比较并交换）来保证事务性操作的逻辑安全。CAS循环的使用方式如下：
\begin{minted}{cpp}
std::atomic<int> counter(0); // 原子变量
// 线程A
int expected = counter.load(); // 读取当前值
int desired = expected + 1; // 计算期望值
while (!counter.compare_exchange_weak(expected, desired)) { // 尝试更新，如果失败，expected会被更新为当前值
    desired = expected + 1; // 重新计算期望值
}
\end{minted}
除了 \verb|compare_exchange_weak|，还有 \verb|compare_exchange_strong|，它们的区别在于前者可能会因为虚假失败而返回 false，而后者则不会。一般来说，如果在循环中使用，推荐使用前者，因为它可能会更快。

此外，原子变量还解决了CPU的OoO问题。OoO指的是CPU为了提高性能，允许某一条指令的操作数好了就提前开跑，而不管这条指令和前后指令的逻辑先后顺序；表现在多线程上，就是一个线程的操作可能会被重排到另一个线程的操作之前，从而导致加了锁也没用。为了解决这个问题，C++提供了内存顺序（memory order）来保证原子变量的操作顺序，能规避OoO问题，有以下几种：
\begin{itemize}
    \item \verb|memory_order_seq_cst|：默认的顺序，保证所有线程看到的操作顺序一致；
    \item \verb|memory_order_acquire|：获取操作，保证在此操作之前的所有读写操作都已经完成，常用于数据的“消费者”获取数据；
    \item \verb|memory_order_release|：释放操作，保证在此操作之后的所有读写操作都不会被重排到此操作之前，常用于数据的“生产者”发布数据；
    \item \verb|memory_order_relaxed|：放宽顺序，不保证操作顺序一致，但可以提高性能，实际是直面OoO，只保证原子操作的原子性，因此建议仅在不涉及数据依赖的情况下使用。
\end{itemize}

\begin{minted}{cpp}
std::atomic<bool> ready{false};
int data = 0;

// 线程甲
data = 42;       
ready.store(true, std::memory_order_release); // 发布（Release）

// 线程乙
while (!ready.load(std::memory_order_acquire)); // 获取（Acquire）
std::println("Data: {}", data); // 读取到的 data 是 42
\end{minted}

这个操作需要注意的另一点是：操作顺序不等于到达顺序。例如以下情景：
\begin{minted}{cpp}
std::atomic<bool> x{false}, y{false};

// 线程甲
x.store(true, std::memory_order_*);
y.store(true, std::memory_order_*);

// 线程乙
std::println("x: {}, y: {}", x.load(std::memory_order_*), y.load(std::memory_order_*));
\end{minted}
读到的 \verb|y| 是true，能否保证读到的 \verb|x| 也是true呢？答案是分情况讨论：
\begin{itemize}
    \item 如果是 \verb|seq_cst|，那可以保证，因为它保证了所有线程看到的操作顺序一致；
    \item 其他：不能保证。它们最多只能保证本线程的指令不重排，不保证另一个线程能感知到本线程内部两个不同变量的先后顺序。
\end{itemize}
因此，如果对两个变量的操作顺序有强依赖，最省事的是全用 \verb|seq_cst|，或者用 \verb|release| 和 \verb|acquire| 来保证数据的“生产者-消费者”关系，在或者使用内存栅栏（memory fence）来保证操作顺序。

\paragraph{信号量} 信号量可以理解为“停车场空位可以有多少”。其内部维护一个计数器，支持三个操作：
\begin{itemize}
    \item \verb|acquire()|：尝试将信号量的计数器减一，如果计数器为零，则阻塞当前线程，直到有其他线程释放信号量，即 \verb|P()| 操作，是原子的；
    \item \verb|release(n=1)|：释放一个信号量，增加计数器的值，并唤醒n个等待的线程（如果有的话，默认为1），即 \verb|V()| 操作，是原子的；
    \item \verb|try_acquire()|：尝试获取信号量，如果计数器为零，则立即返回 false，否则减一并返回 true。
\end{itemize}
信号量是没有所有权的，任何线程都可以释放信号量。这实际上是很灵活的，但也很危险。如果某个线程因为异常或逻辑错误少调用了 \verb|release()|，那么其他线程就可能永远阻塞，导致假死。此外，信号量有硬上限，可以传入一个较大的数值。此外，在极端的高性能场景，其调用开销不如无锁队列；但对于常规业务已经基本足够。最后，信号量不保护，只计数，因此要保护的话还是要加锁。

另一个信号量是 \verb|std::binary_semaphore|，它是一个二值信号量，计数器只能为0或1。强烈不建议用它替代 \verb|std::mutex|。

其典型用途是限流，比如限制同时访问某个资源的线程数，或者实现生产者-消费者模型。

\begin{minted}{cpp}
#include <semaphore>

std::counting_semaphore<3> sem(3); // 初始化信号量，计数器为3

void heavy_task(int id) {
    sem.acquire(); // 尝试获取信号量，如果计数器为0，则阻塞
    std::cout << "Thread " << id << " is doing heavy task." << std::endl;
    std::this_thread::sleep_for(std::chrono::seconds(2)); // 模拟耗时任务
    std::cout << "Thread " << id << " finished heavy task." << std::endl;
    sem.release(); // 释放信号量，增加计数器并唤醒一个等待的线程
}

int main() {
    std::vector<std::thread> threads;
    for (int i = 0; i < 10; ++i) {
        threads.emplace_back(heavy_task, i); // 创建10个线程
    }
    for (auto& t : threads) {
        t.join(); // 等待所有线程完成
    }
    return 0;
}
\end{minted}
这个代码中，控制台会每隔2秒输出3行，表示同时只有三个线程在执行，而其他线程则在等待信号量的释放。

或者把它当成纯计数器也行。它实现了和原子量、条件变量类似的功能：当队列中没有东西的时候，消费者会睡觉、等着生产者产生新的数据，而不是空转CPU。
\begin{minted}{cpp}
#include <queue>
#include <mutex>
#include <semaphore>

std::queue<int> q;
std::mutex mtx;
std::counting_semaphore<1000> data_count(0); // 假设队列不会这么长

// 生产者
void producer() {
    for (int i = 0; i < 10; ++i) {
        {
            std::lock_guard<std::mutex> lock(mtx);
            q.push(i);
        }
        data_count.release(); // 通知有一个数据可用
        std::cout << "Produced: " << i << std::endl;
    }
}

// 消费者
void consumer() {
    for (int i = 0; i < 10; ++i) {
        data_count.acquire(); // 等待数据（计数 > 0）
        std::lock_guard<std::mutex> lock(mtx);
        int val = q.front();
        q.pop();
        std::cout << "Consumed: " << val << std::endl;
    }
}
\end{minted}

\paragraph{任务打包} 任务打包（\verb|std::packaged_task|）是一个模板类，它允许你将一个可调用对象（函数、函数指针、lambda表达式等）包装成一个任务，并与一个 \verb|std::future| 对象关联，从而可以在未来的某个时间点获取任务的结果。任务打包的使用方式如下：
\begin{minted}{cpp}
std::packaged_task<int(int, int)> task([](int a, int b) { return a + b; }); // 打包一个加法任务
std::future<int> result = task.get_future(); // 获取与任务关联的 future
task(2, 3); // 执行任务
std::println("Result: {}", result.get()); // 获取结果并打印
\end{minted}
这个东西不能复制，只能移动，只能调用一次，且会把异常传进 \verb|std::future| 对象中。当调用 \verb|result.get()| 时，如果任务抛出了异常，那么这个异常会被重新抛出，从而可以在调用方捕获。

\paragraph{栅栏} 可复用栅栏（\verb|std::barrier|）和一次性栅栏（\verb|std::latch|）是两个同步原语，它们都用于线程之间的同步，但有不同的使用场景。可复用栅栏是阶段性的同步，而一次性栅栏是一次性的同步。注意：这两个栅栏不是内存栅栏（fence）。

一次性栅栏是一个向下的计数器，在创建时就设定一个初始值。不同的线程可以递减这个计数器或阻塞等待，直到计数器归零。当计数器归零后，该栅栏被释放，所有等待的线程都会被唤醒。一次性栅栏的使用方式如下：
\begin{minted}{cpp}
#include <latch>
std::latch latch(3); // 创建一个计数器为3的栅栏
// 子线程们
{
    // working...
    latch.count_down(); // 递减计数器
}

// 主要线程
{
    latch.wait(); // 阻塞等待，直到计数器归零
    // 所有子线程都完成了工作，可以继续执行
}
\end{minted}
一次性栅栏是无锁的，适合用于一次性的同步场景。

可复用栅栏与前者类似，但它可以在每个阶段结束后被重置，从而可以在多个阶段中使用。可复用栅栏的使用方式如下：
\begin{minted}{cpp}
#include <barrier>
std::barrier barrier(3); // 创建一个计数器为3的栅栏
// 子线程们
{
    for (int phase = 0; phase < 3; ++phase) {
        // 模拟一个多阶段且需要常常同步的工作
        barrier.arrive_and_wait(); // 递减计数器并阻塞等待，直到计数器归零
        // 所有线程都到达了栅栏，可以继续执行
    }
}
// 主要线程
{
    for(auto& t : threads) {
        t.join(); // 等待所有子线程完成
    }
}
\end{minted}
注意：构造栅栏时指定的线程数量必须和实际参与的线程数目一致。如果有一个线程需要中途退出，那么就必须在退出前调用 \verb|barrier.arrive_and_drop()|，以减少后续阶段的参与计数，否则会导致其他线程一直阻塞等待，最终导致假死。

\subsection{竞态条件、死锁和逻辑死锁}

虽然在RAII和线程池的设计下把这三尊大佛规避了大半，但实际情形是如果逻辑错误，那这些条件都有可能发生。下面就拿出来讲讲怎么预防它们。

\paragraph{竞态条件}

在多个线程同时访问共享资源时，如果有多个线程试图读写该资源而没有明显的顺序，就可能会出现竞态条件（Race Condition）。竞态条件发生的主要原因是没有加锁或没有按原子操作，导致多个线程同时修改共享资源，从而产生不可预测的结果。竞态条件的典型结果是：程序能编译，能执行完，但结果不符合预期，且多次执行得到的结果不同。

比如说下面这个代码就会产生竞态条件：
\begin{minted}{cpp}
#include <iostream>
int counter = 0; // 共享资源
void increment() {
    for (int i = 0; i < 100000; ++i) {
        ++counter; // 竞态条件：多个线程同时修改 counter
    }
}
int main() {
    std::thread t1(increment);
    std::thread t2(increment);
    t1.join();
    t2.join();
    std::cout << "Counter: " << counter << std::endl; // 结果可能不是 200000
    return 0;
}
\end{minted}
为了避免竞态条件，我们有三个主要手段：
\begin{itemize}
    \item 上锁：
    \begin{minted}{cpp}
std::mutex mtx; // 互斥量
void increment() {
    for (int i = 0; i < 100000; ++i) {
        std::lock_guard<std::mutex> lock(mtx); // 上锁
        ++counter; // 访问共享资源
    }
}
    \end{minted}
    \item 原子操作：
    \begin{minted}{cpp}
std::atomic<int> counter(0); // 原子变量
void increment() {
    for (int i = 0; i < 100000; ++i) {
        ++counter; // 原子操作，线程安全
    }
}
    \end{minted}
    \item 传拷贝而不是引用，避免共享状态。
\end{itemize}

\paragraph{死锁}

死锁是指多个线程在等待彼此持有的资源，从而导致所有线程都无法继续执行。死锁发生的主要原因是：
\begin{itemize}
    \item 不可剥夺：线程持有的资源不能被其他线程强行夺走；
    \item 占有且等待：线程自己持有资源的同时又在等待其他线程持有的资源；
    \item 互斥：线程对资源的访问是互斥的，即同一时间只能有一个线程访问资源；
    \item 循环等待：线程形成一个环形的等待链，每个线程都在等待下一个线程持有的资源。
\end{itemize}
就相当于：车钥匙在家里，家门钥匙却在车里，而且车钥匙和家门钥匙都只有一把，导致车和家都无法进入。死锁的主要表现是：程序卡住了，无法继续执行，且没有任何错误提示，且并不会很吃CPU资源。

只要废掉这四个条件中的任何一个，就可以避免死锁的发生：
\begin{itemize}
    \item 废掉不可剥夺：可以使用信号量或条件变量来实现资源的可剥夺性，即当一个线程持有资源时，如果其他线程需要该资源，可以强制该线程释放资源，从而避免死锁。举例：请开锁师傅帮忙开其中一个门。
    \item 废掉占有且等待：可以要求线程在获取资源时一次性获取所有需要的资源，或者在获取资源时不持有其他资源，从而避免死锁。举例：车钥匙和家门钥匙都串在一起，一次性拿走。
    \item 废掉互斥：可以使用读写锁或共享锁来允许多个线程同时访问资源，从而避免死锁。举例：准备备用钥匙。
    \item 废掉循环等待：可以给资源分配一个全局的顺序号，使得线程在请求资源时按照顺序进行，从而避免循环等待。举例：一开始就别把房门钥匙锁到车里。实际上这是破除死锁的核心法则。
\end{itemize}

实际情况下一般按下面四步骤操作：
\begin{enumerate}
    \item 固定锁顺序，约定所有线程都必须按相同的顺序加锁。只要顺序一致就不可能死锁。
    \item 一次性获取所有锁，避免占有且等待。
    \item 少加锁，只有真正操作共享数据的临界区内才上锁，否则坚决不加锁。
    \item 如果长期拿不到锁，就放弃当前操作，主动释放已有的锁，防止无限僵持。
\end{enumerate}

\paragraph{逻辑死锁}

逻辑死锁是指程序在逻辑上陷入了死循环，但不是因为资源竞争导致的死锁。比如说，信号量丢失（少 \verb|release()| 了一次导致之后的线程等待一次永远都不会来的 \verb|release()|）、参与量计数永久未达标（少 \verb|arrive_and_drop()| 了一次导致之后的线程等待一次永远都不会来的同步）等。这种情况是有线程等待一个永远不来的事件，虽然不符合死锁的经典条件，效果却一致：整个系统会卡死。

对于这种情况，最好的办法就是：RAII，或永远不要绕过计数器直接操作。这个和我们先前讨论的抛出异常时对资源的处理方法类似：要么RAII，要么在可能抛异常的地方之后再开始计数，再或者在即将抛异常之前释放计数。总之，不RAII的情况下，信号量或参与量必须被正确的调用。

比如我们来RAII地包装一个信号量：
\begin{minted}{cpp}
template <int T>
class CountingSemaphoreGuard{
    std::counting_semaphore<T>& sem;
public:
    CountingSemaphoreGuard(std::counting_semaphore<T>& s) : sem(s) { sem.acquire(); } // 构造时获取信号量
    ~CountingSemaphoreGuard() { sem.release(); } // 析构时释放信号量
    CountingSemaphoreGuard(const CountingSemaphoreGuard&) = delete; // 禁止拷贝
    CountingSemaphoreGuard& operator=(const CountingSemaphoreGuard&) = delete; // 禁止赋值
    CountingSemaphoreGuard(CountingSemaphoreGuard&&) = delete; // 禁止移动
}; // 之后就用它来RAII的包装信号量，保证每次 acquire() 都有对应的 release()，从而避免逻辑死锁。
\end{minted}
极端情况下，也可以用超时重置机制。但这是最后的手段。

\section{无锁数据结构}

刚刚我们写的线程池，每次对 \verb|tasks| 队列的访问都需要加锁。加锁和解锁的开销不算小，因为有CPU的上下文切换开销。那么有没有一种办法，能彻底规避锁的开销？这就是无锁数据结构。

\subsection{无锁数据结构的原理}

无锁数据结构的核心就是\textbf{原子操作}，因为原子操作是CPU提供的最小粒度的同步原语，它可以保证在多线程环境下对共享数据的访问不会发生竞态条件。

以无锁队列为例，其工作逻辑是这样的：
\begin{itemize}
    \item 读快照：读当前的队列头和队列尾的指针，形成一个快照；
    \item 计算新值：根据快照计算出新的队列头和队列尾的指针；
    \item 原子试探：使用CAS（Compare-And-Swap）操作尝试将队列头和队列尾的指针更新为新的值。
    \item 重试循环：如果CAS操作失败，说明有其他线程修改了队列头和队列尾的指针，那么就重新读快照，计算新值，再尝试原子更新，直到成功为止。这会在一个循环（“自旋”）中不断尝试，线程不会休眠，直到成功为止。
\end{itemize}

因此，无锁队列的工作逻辑是这样的：
\begin{minted}{cpp}
#include <atomic>
template<typename T>
class LockFreeQueue {
    struct Node {
        T data;
        std::atomic<Node*> next;
        Node(T val) : data(val), next(nullptr) {}
    };
    std::atomic<Node*> head;
    std::atomic<Node*> tail;
public:
    LockFreeQueue() {
        Node* dummy = new Node(T()); // 创建一个哑节点
        head.store(dummy);
        tail.store(dummy);
    }
    ~LockFreeQueue() {
        while (Node* node = head.load()) {
            head.store(node->next.load());
            delete node;
        }
    }
    void enqueue(T val) {
        Node* newNode = new Node(val);
        Node* oldTail;
        do {
            oldTail = tail.load();
            Node* next = oldTail->next.load();
            if (next == nullptr) { // 如果尾节点的 next 为空，说明可以插入
                if (oldTail->next.compare_exchange_weak(next, newNode)) {
                    break; // 插入成功，跳出循环
                }
            } else { // 尾节点的 next 不为空，说明有其他线程在插入，尝试移动尾指针
                tail.compare_exchange_weak(oldTail, next);
            }
        } while (true);
        tail.compare_exchange_weak(oldTail, newNode); // 尝试移动尾指针到新节点
    }
    bool dequeue(T& val) {
        Node* oldHead;
        do {
            oldHead = head.load();
            Node* oldTail = tail.load();
            Node* next = oldHead->next.load();
            if (oldHead == oldTail) { // 队列为空或只有一个节点
                if (next == nullptr) {
                    return false; // 队列为空
                }
                tail.compare_exchange_weak(oldTail, next); // 尝试移动尾指针
            } else {
                val = next->data; // 获取下一个节点的数据
                if (head.compare_exchange_weak(oldHead, next)) {
                    delete oldHead; // 删除旧的头节点
                    return true; // 出队成功
                }
            }
        } while (true);
    }
};

LockFreeQueue<std::function<void()>> tasks; // 无锁队列

void DenseTask(int a) {
    std::this_thread::sleep_for(std::chrono::seconds(1));
    std::cout << "Task " << a << " done." << std::endl;
}

for (int i = 0; i < 8; ++i) {
    tasks.enqueue([i]() { DenseTask(i); }); // 提交任务到无锁队列
}
\end{minted}

无锁（自旋锁）和互斥锁（阻塞锁）的区别见表~\cref{tab:lock}。

\begin{table}
    \centering
    \begin{tabular}{l|p{6cm}p{6cm}}
        \toprule
        & 互斥锁 & 自旋锁 \\
        \midrule
        锁的本质 & 阻塞其他线程 & 原子操作 \\
        线程挂起 & 拿不到锁就睡、让出CPU（消极等待） & 拿不到锁就自旋，霸占CPU（积极等待）\\
        延迟 & 有上下文切换（微秒级） & 几乎没有 \\
        死锁风险 & 有 & 没有锁的持有关系，不存在死锁 \\
        活锁风险 & 不存在 & 两个线程可能互相谦让而导致长期失败 \\
        内存回收 & 容易（析构时释放锁） & 复杂，需要保证没有线程在自旋时访问已释放的内存 \\
        适用场景 & 绝大多数情况 & 极致性能，或绝对不能休眠的中断上下文和内核态 \\
        \bottomrule
    \end{tabular}
    \caption{互斥锁和自旋锁的区别}
    \label{tab:lock}
\end{table}

\subsection{无锁数据结构的致命问题} 

虽然无锁数据结构比正常的有锁数据结构更快，但它们也有一些致命的问题，主要包括ABA问题、内存回收问题和活锁问题。这和有锁数据结构不同：有锁数据结构只需要维护锁的持有关系不乱，且防止OoO对多线程的影响就可以了，但无锁数据结构的这些问题解决难度更大。

\paragraph{ABA问题} ABA问题的现象是这样的：线程甲读取一个栈顶为A，准备把A换成B。正在计算B的时候，线程乙把A给删了，又新建了一个地址和A相同的C压回去。线程甲回来时，通过CAS判断A和C是一个东西，于是把指针改成了B，但这个时候 \verb|B->next| 是已经被删除的旧A，于是导致崩溃。

这个问题的核心原因是：线程分不清A和C是不是同一个东西。为了解决这个问题，必须引入\textbf{带版本号的指针}，每次修改的时候都让版本号变化。这样，即使地址相同，版本号变了也能识别出来。

\paragraph{内存回收问题} 这个问题的现象是这样的：线程甲读取某节点的数据，线程乙却CAS成功了，把这个节点删除了并释放了内存。线程甲继续访问这个节点的数据，结果它的指针就成了野指针，引发未定义行为。

这个问题的核心原因是线程不知道有没有其他线程在访问现有的节点。解决方法有两种：要么引入“风险指针”，让其他的线程知道有东西在读这个，不让它删；要么引入基于时代的回收策略，等所有线程都进入下一个“时代”后再删除那些该删除的节点。

\paragraph{活锁} 活锁是CAS中最常见的问题。当两个线程都在自旋、读取同一个变量时，它们可能会互相干扰，对方的修改总是覆盖自己的修改，导致永远无法完成自己的目标，即互相谦让导致的长期失败。如果说死锁是两辆车在窄巷相遇、谁也不愿倒车，那活锁就是两辆车都极有礼貌的倒出巷子再回来，结果在巷子中又碰面了，但两位司机都有点傻，不会等到另一辆车从巷子中出来后再继续行驶，如此形成无限循环。由此可见，这是自旋锁特有的风险，互斥锁不可能出现这种情况，因为互斥锁是阻塞的，拿不到锁就睡觉、让出CPU。

活锁的杀伤力比死锁大得多。死锁可以通过加锁的顺序规避，即使在编程中未能规避，发生死锁时CPU立即停摆，程序卡住，监控系统会告警，运维人员可以迅速介入，也容易察觉发生死锁的位置。但活锁会伪装：不仅很难在代码中规避，而且发生时CPU仍然在疯狂自旋，QPS（每秒查询率）也非常高，监控系统也不会告警，唯一的线索是业务数据库中的关键数据就是不更新，不容易意识到是活锁的问题，而且排查时也很难定位到问题的根源。

解决活锁有三种手段：
\begin{description}
    \item[引入随机退避机制] 线程失败后，随机休眠几个微秒或自旋几次后再尝试，错开“撞车”时间。这是Ethernet的CSMA/CD协议的核心思想。
    \item[指数退避] 线程失败后，休眠一定时间，如果还失败则休眠更长时间，直到成功为止。这个做法类似于TCP的拥塞控制算法，也是工业级别无锁队列的常用做法。
    \item[让出时间片] 如果是在同一个核心上超线程竞争，可以让出时间片，给其他线程机会。
\end{description}

实际情况中，活锁的发生概率和严重程度取决于线程数、CPU核心数、线程调度策略、线程优先级等因素。一般来说，线程数越多、CPU核心数越少、线程调度策略越公平、线程优先级越接近，活锁的发生概率就越高，严重程度也就越大。

\subsection{无锁数据结构的应用}

无锁数据结构虽然够快，但我非常不建议使用无锁的数据结构，尤其是手搓。这是因为无锁数据结构非常容易出错，而且一旦出错，排查起来非常困难。除非你是一个非常有经验的并发程序员，否则不要轻易尝试，而是改用有锁数据结构代替。因为有锁数据结构对大多数场景而言确实已经足够快，且代码正确性容易证明；此外，无锁带来的收益在IO密集型场景下微乎其微。

真的硬要用，那建议使用成熟的无锁数据结构库，包括：
\begin{itemize}
    \item \verb|moodycamel::ConcurrentQueue|：高吞吐队列，工业级优化。
    \item \verb|boost::lockfree::queue|：Boost库提供的简单栈和队列。
\end{itemize}
最终依然是那句话：先保证对，再保证快；能用锁解决就不要玩无锁，要用就用现成的，不要自己手搓。